\documentclass[10pt]{article}
\usepackage[T1]{fontenc}
\usepackage[utf8]{inputenc}
\usepackage[french]{babel}
\usepackage{amsmath,amssymb,mathrsfs,mathpazo}
\usepackage{geometry}
\geometry{a4paper,margin=2cm}
\usepackage{graphicx,array,tabularx,tabularray,colortbl,multirow,multicol,booktabs,enumitem,pifont,ulem,fancybox,diagbox,cancel}
\usepackage{pstricks,pstricks-add,pst-plot,pst-tree,pst-node,pst-text,pst-func,pst-math,pst-bspline,pst-3dplot,pst-eucl,pst-solides3d}
\let\pscalculate\undefined
\usepackage{pst-bezier}
\usepackage{tikz,pgfplots,tkz-tab}
\usetikzlibrary{arrows,arrows.meta,patterns,calc,positioning,trees,decorations.markings,shapes.geometric,shapes.misc,shapes.arrows,intersections,decorations.pathmorphing,decorations.pathreplacing,backgrounds,shadows,babel}
\pgfplotsset{compat=1.18}
\usepgfplotslibrary{fillbetween}
\usepackage[np]{numprint}
\usepackage{listings,makecell,pas-tableur,twoopt}
\let\barre\undefined
\usepackage[pstricks]{tablvar}
\definecolor{darkgreen}{RGB}{0,100,0}
\providecommand{\red}{\color{red}}
\providecommand{\blue}{\color{blue}}
\providecommand{\green}{\color{green}}
\providecommand{\black}{\color{black}}
\providecommand{\white}{\color{white}}
\providecommand{\gray}{\color{gray}}
\providecommand{\cyan}{\color{cyan}}
\providecommand{\magenta}{\color{magenta}}
\providecommand{\yellow}{\color{yellow}}
\providecommand{\violet}{\color{violet}}
\providecommand{\orange}{\color{orange}}
\providecommand{\euro}{\texteuro}
\providecommand{\EUR}{\texteuro}
\providecommand{\ieme}{\textsuperscript{e}}
\providecommand{\degres}{\ensuremath{^\circ}}
\providecommand{\degree}{\ensuremath{^\circ}}
\providecommand{\vv}[1]{\overrightarrow{#1}}
\providecommand{\Coord}[1]{#1}
\providecommand{\decosix}{}
\providecommand{\starredbullet}{\ensuremath{\bullet}}
\pagestyle{empty}
\setlength{\parindent}{0pt}
\usepackage{adjustbox}\DeclareUnicodeCharacter{2236}{:}\DeclareUnicodeCharacter{2190}{\ensuremath{\leftarrow}}
\providecommand{\R}{\mathbb{R}}
\providecommand{\N}{\mathbb{N}}
\providecommand{\Z}{\mathbb{Z}}
\providecommand{\Q}{\mathbb{Q}}
\providecommand{\D}{\mathbb{D}}
\providecommand{\C}{\mathbb{C}}
\providecommand{\e}{{\mathrm{e}}}
\providecommand{\vect}[1]{\overrightarrow{#1}}
\providecommand{\vectt}[1]{\overrightarrow{\mathrm{#1}}}
\providecommand{\vvt}[1]{\overrightarrow{\mathrm{#1}}}
\providecommand{\np}[2][]{{#2~#1}}
\providecommand{\vv}[1]{{\overrightarrow{#1}}}
\renewcommand{\d}{{\mathrm{d}}}
\providecommand{\degres}{{^\circ}}
\providecommand{\degre}{{^\circ}}
\providecommand{\blue}{\color{blue}}
\providecommand{\red}{\color{red}}
\providecommand{\enskip}{\quad}
\providecommand{\llbracket}{\lbrack\!\lbrack}
\providecommand{\rrbracket}{\rbrack\!\rbrack}
\providecommand{\euro}{{\text{€}}}
\providecommand{\no}{{\text{n°}}}
\providecommand{\interval}[2]{{[#1;#2]}}
\providecommand{\intervFO}[2]{{[#1;#2[}}
\providecommand{\dotfill}{\ldots}

\providecommand{\ds}{\displaystyle}
\providecommand{\textding}[1]{#1}
\providecommand{\Oij}{\(\mathrm{O};\vec{\imath},\vec{\jmath}\)}
\providecommand{\Oijk}{\(\mathrm{O};\vec{\imath},\vec{\jmath},\vec{k}\)}
\providecommand{\Ouv}{\(\mathrm{O};\vec{u},\vec{v}\)}
\providecommand{\intervFF}[2]{\ensuremath{[#1;#2]}}
\renewcommand{\euro}{\ifmmode\text{\texteuro}\else\texteuro\fi}
\DeclareUnicodeCharacter{00B0}{\ifmmode^{\circ}\else\textdegree\fi}
\DeclareUnicodeCharacter{20AC}{\ifmmode\text{\texteuro}\else\texteuro\fi}
% Énoncés : Éducation nationale. Transcription et corrigés : APMEP, Denis Vergès et contributeurs. Découpage et classement : CoopMaths, Sébastien Lozano, Éric Elter et contributeurs. Adaptation de la présentation : Maths Explorer.
% CC BY-SA : https://coopmaths.fr/www/about/
% https://coopmaths.fr/alea/static/eam/2026/tex/eam_technologique_2026_09_metropole_automatismes_cor.tex
\begin{document}
\noindent\textit{APMEP / CoopMaths. CC BY-SA. Présentation adaptée par Maths Explorer.}\par\medskip
\textbf{Pour cette première partie, aucune justification n'est demandée et une seule réponse est possible par question.} Les justifications ci-dessous sont fournies à titre pédagogique.

Une réponse fausse ou l'absence de réponse n'enlève aucun point : il ne faut donc jamais laisser une question sans réponse.

\subsection*{Question 1 \quad Réponse A}

On développe le carré avec $(a-b)^2 = a^2 - 2ab + b^2$, puis on réduit :
\[f(x) = (x-1)^2 - 8 = x^2 - 2x + 1^2 - 8 = x^2 - 2x - 7.\]

\subsection*{Question 2 \quad Réponse D}

Une baisse de 20\,\% correspond au coefficient multiplicateur $0,8$ ; une baisse de
10\,\%, au coefficient multiplicateur $0,9$. Le coefficient global est le produit :
\[0,8 \times 0,9 = 0,72 = 1 - 0,28.\]
Le prix a donc baissé de 28\,\%.

\subsection*{Question 3 \quad Réponse A}

Les solutions de $f(x) = 0$ sont les abscisses des points où $\mathcal{C}_f$ coupe l'axe des
abscisses. La courbe le coupe en trois points, d'abscisses $-5$, $1$ et $6$. L'équation a
donc trois solutions, dont deux positives (1 et 6).

\subsection*{Question 4 \quad Réponse D}

La droite $(d)$ coupe l'axe des ordonnées en 3 : $b = 3$. Elle passe aussi par le point
de coordonnées $(6~;~0)$, donc
\[a = \dfrac{0 - 3}{6 - 0} = -\dfrac{3}{6} = -\dfrac{1}{2}.\]
L'équation réduite est $y = -\dfrac{1}{2}x + 3$.

\subsection*{Question 5 \quad Réponse B}

Dans chaque barre, les nuitées en camping correspondent à la hauteur de la partie
noire (et non à la hauteur de son sommet). Par lecture graphique, cette partie mesure
environ 80~millions de nuitées en 2016 (de 200 à 280 environ), alors qu'elle dépasse
100~millions toutes les autres années.

\subsection*{Question 6 \quad Réponse C}

Multiplier par $10^{-3}$ revient à décaler la virgule de trois rangs vers la gauche :
$N = \np{0,000026}$. En écriture scientifique :
\[N = 0,026 \times 10^{-3} = 2,6 \times 10^{-2} \times 10^{-3} = 2,6 \times 10^{-5}.\]

\subsection*{Question 7 \quad Réponse A}

Le coefficient directeur de la droite (AB) est
\[a = \dfrac{y_B - y_A}{x_B - x_A} = \dfrac{-5 - 15}{7 - 5} = \dfrac{-20}{2} = -10.\]

\subsection*{Question 8 \quad Réponse D}

$4x + 3 = 0 \iff x = -\dfrac{3}{4}$. Le coefficient de $x$ vaut $4 > 0$ : la fonction
affine est croissante, donc négative avant $-\dfrac{3}{4}$ et positive après.

\subsection*{Question 9 \quad Réponse A}

On réduit d'abord le dénominateur, puis on prend l'inverse :
\[2 + \dfrac{5}{6} = \dfrac{12}{6} + \dfrac{5}{6} = \dfrac{17}{6},
\qquad\text{donc}\qquad
F = \dfrac{1}{\frac{17}{6}} = \dfrac{6}{17}.\]

\subsection*{Question 10 \quad Réponse A}

Un produit de facteurs est nul si et seulement si l'un au moins de ses facteurs est nul :
\[(3x+1)(x-5) = 0 \iff 3x + 1 = 0 \ \text{ ou } \ x - 5 = 0
\iff x = -\dfrac{1}{3} \ \text{ ou } \ x = 5.\]
D'où $S = \left\{-\dfrac{1}{3}~;~5\right\}$.

\subsection*{Question 11 \quad Réponse D}

Les solutions de $f(x) \leqslant g(x)$ sont les abscisses des points où $\mathcal{C}_f$ est
sur ou sous $\mathcal{C}_g$. Les deux courbes se coupent aux points de coordonnées
$(-3~;~-1)$ et $(2~;~3)$ ; la parabole $\mathcal{C}_f$ est au-dessus de la droite
$\mathcal{C}_g$ entre $-3$ et $2$, et en dessous ailleurs sur $[-5~;~5]$. D'où
$S = [-5~;~-3] \cup [2~;~5]$.

\subsection*{Question 12 \quad Réponse B}

Par lecture graphique, $\mathcal{C}_f$ part du point de coordonnées $(-5~;~0)$, monte jusqu'au point $(-2~;~5)$, descend jusqu'au point $(4~;~6)$, puis remonte jusqu'au point $(6~;~0)$. Dans un tableau de variations, la première ligne contient les abscisses $-5$, $-2$, $4$, $6$,
et les flèches relient les ordonnées $0$, $5$, $-6$, $0$ : c'est le tableau B.

\end{document}
